\section{Computational demonstration: discontinuity in $\mathbb{C}$ and a continuous lift in $\mathbb{R}^4$}

\subsection{Aim}
The conjecture asserts that inverse hyperoperators may require additional
``information storage'' beyond the ambient value space. In the concrete
test case $f(z)=z^z$, the inverse problem
\[
z^z = w
\]
is multivalued. Our computational aim is to exhibit \emph{monodromy} of the inverse:
there exist closed loops $\gamma:[0,1]\to\Omega\subset\mathbb{C}$ in the $w$-plane
such that continuing a local inverse branch along $\gamma$ returns to a different
branch at the same basepoint.

Equivalently, if one attempts to represent the inverse as a single-valued
continuous function $s:\Omega\to\mathbb{C}$, then tracking $s(\gamma(t))$ must
produce a discontinuity for suitable loops.

\subsection{What the computation does}
We choose two explicit loops in the $w$-plane:
\[
\gamma_{0}(t) = w_0 + R e^{2\pi i t} \quad \text{(a loop around the origin)}
\]
and
\[
\gamma_{w^\ast}(t) = w_1 + R e^{2\pi i t} \quad \text{(a loop around a critical value $w^\ast$).}
\]
For each loop $\gamma(t)$ we numerically track a continuous choice of preimage
$z(t)$ satisfying $z(t)^{z(t)} = \gamma(t)$ by nearest-neighbour continuation
among nearby candidate branches.

We then measure the \emph{gap} in $\mathbb{C}$:
\[
\mathrm{gap} = |z(1)-z(0)|,
\]
and the step-size spike in the tracked path, which signals a branch jump
(discontinuity in any attempted single-valued complex representation).

\subsection{Lift coordinates (no interpolation)}
To model the ``extra information'' carried by branching, we introduce two
real lift coordinates computed directly from the path (not interpolated):
\[
\theta(t) = \mathrm{unwrap}\,\arg(\gamma(t)),
\qquad
c(t) = \log|\gamma(t)| + i\theta(t),
\qquad
\phi(t) = \mathrm{unwrap}\,\arg\big(c(t)+1/e\big),
\]
where $c=-1/e$ is the principal branch point of the Lambert $W$ function.
We then define a lifted path in $\mathbb{R}^4$ by
\[
Z(t) = \big(\Re z(t),\, \Im z(t),\, \theta(t)/(2\pi),\, \phi(t)/(2\pi)\big).
\]

\subsection{What the computation shows (and what it does not)}
The computation shows:
\begin{itemize}
\item The tracked inverse path is typically discontinuous in $\mathbb{C}$
for loops with nontrivial monodromy (nonzero $\mathrm{gap}$ and a large
step-size spike).
\item The lifted path $Z(t)\in\mathbb{R}^4$ is continuous (step sizes remain bounded),
and the endpoint displacement in the lift coordinates records sheet change.
\end{itemize}

Importantly, on a covering-space lift a loop in $\Omega$ need not close upstairs:
the lifted path may end on a different sheet (as in the standard logarithm
Riemann surface). The claim here is continuity of the lift, not closure.

This provides computational evidence for the ``information storage'' perspective:
two independent sources of multivaluedness require two additional real lift
coordinates beyond the complex value coordinates.
